\section{Formal Analysis}\label{sec:formal}

We model \sys{}, four baselines, a mutant and a straw man in Tamarin 1.10.0~\citep{meier2013tamarin}\footnote{Tamarin
reports the Maude 3.2 rewriting engine we used as unsupported (it recommends 3.2.1 or later); all verdicts were
obtained with 3.2.} to check that the binding properties
hold for an unbounded number of sessions, and to make the assumptions behind them explicit.
Formal models of agent authorization protocols already exist, including Tamarin analyses of agent payment
protocols~\citep{jiang2026payment} and of a post-quantum delegation protocol~\citep{chen2026aith}; we claim no
priority for the method.

\subsection{Model and hypotheses}

The agent is a Dolev-Yao adversary. It chooses every tool call, including the tool, target and amount of
any effect and any explanation text, sees everything the gateway
returns and every WebAuthn assertion, and can send any message to any gateway rule. The challenge is modeled as
the free tuple
\[\langle\texttt{'ASU-APPROVAL-v2'},\ \dig,\ \mathit{id},\ g\rangle,\]
This is sound because the canonical
JSON encoding of an object with a fixed key set is injective and SHA-256 is assumed collision resistant. The
nonce is dropped, which only makes signatures reusable in more contexts and is therefore conservative for the
safety lemmas. Rendering is injective, which the code discharges with a property test.

Three hypotheses are stated in the theory, not hidden in it. Each is a rule the adversary may use, each rule
emits a marker, and each lemma names the hypotheses it relies on as conjuncts of its premise:
\begin{itemize}[leftmargin=1.4em,itemsep=1pt]
\item[\Hyp{1}] The human decides from the trusted fields of the card, not from the agent's explanation.
\item[\Hyp{2}] Enrollment is closed: only the device holder can enroll a principal's first credential.
\item[\Hyp{3}] The rendering origin is honest: the page that drives the authenticator shows the card that
  matches the challenge it hands over. Authenticators do not display the challenge, so this is a real assumption.
\end{itemize}

\subsection{Results}

All \foLemmas{} lemmas across \foTheories{} theories produced the expected verdict, with \foMismatch{}
mismatches and no errors; the full run takes about three minutes, and the slowest lemma, which bounds the number
of grant uses, takes \foMaxSec{} seconds. \Cref{tab:formal} reports the main lemmas.

\begin{table}[t]
\centering\small
\caption{Main \sys{} lemmas (Tamarin 1.10.0, unbounded). ``Hyp.'' lists the hypotheses in the premise. The
existence lemmas show that executions, grant executions and receipt verification are reachable under
\Hyp{1} to \Hyp{3}, so the safety lemmas are not vacuous.}
\label{tab:formal}
\begin{tabularx}{\textwidth}{@{}l>{\raggedright\arraybackslash}X l l@{}}
\toprule
Lemma & Statement & Hyp. & Result\\
\midrule
\code{L1\_effect\_binding} & An execution of $\Desc$ under $\mathit{id}$ implies the human approved $\Desc$ for $\mathit{id}$ beforehand & \Hyp{1}--\Hyp{3} & verified\\
\code{L1\_sig\_binding} & An execution implies $U$'s authenticator signed $\dig(\Desc)$ for $\mathit{id}$ & \Hyp{2} & verified\\
\code{L2\_single\_use} & At most one execution per approval identifier & none & verified\\
\code{L3\_agent\_independence} & An approved status implies the human approved that effect & \Hyp{1}--\Hyp{3} & verified\\
\code{L4\_non\_repudiation} & A judge holding only $pk_U$ accepts $(\dig,\mathit{id},g)$ only if $U$'s authenticator signed exactly that tuple, even when the origin controls the ceremony & \Hyp{2} & verified\\
\code{L4c\_consent\_nonrep.} & A judge's acceptance implies the human approved the effect & \Hyp{1}--\Hyp{3} & verified\\
\code{L5\_grant\_scope} & A grant execution of $\langle t,r,a\rangle$ implies the human approved $\langle t,r,a_0\rangle$ and the gateway-built grant, with $a\le a_0$ & \Hyp{1}--\Hyp{3} & verified\\
\code{L5b\_grant\_max\_uses} & At most \emph{max\_uses} grant executions & none & verified\\
\midrule
\code{X1\_*} (2 lemmas) & L1 and L3 with \Hyp{1} dropped: the human decides on the explanation & & falsified\\
\code{X2\_*} (2 lemmas) & L4c and L1 with \Hyp{3} dropped: the origin shows one card and has another signed & & falsified\\
\code{X3\_*} (4 lemmas) & L1, L1\_sig, L4 and L5 with \Hyp{2} dropped: the agent enrolls first & & falsified\\
\bottomrule
\end{tabularx}
\end{table}

\paragraph{Each hypothesis is necessary.} The X variants are the same theory as the main lemmas with exactly one
hypothesis removed and the conclusion unchanged, so the attacked model cannot drift from the verified one. Each
variant is falsified, and Tamarin returns the expected attack. Without \Hyp{1}, the agent writes ``send 5 to
alice'' as its explanation for an effect that pays 500 to an attacker; the human approves on the explanation and
taps, yet \code{L1\_sig\_binding} and \code{L4\_non\_repudiation} still hold, because the signature binds the real
digest. Without \Hyp{3}, the origin shows one card and has the authenticator sign another approval's challenge;
the key-level lemma \code{L4\_non\_repudiation} still verifies in the same theory, which is why we state
non-repudiation at the key level (\Prop{5}). Without \Hyp{2}, the agent enrolls the first credential of an unused
principal and approves its own calls, which is \AB{7}.

\paragraph{The baselines fail where they should.} \Cref{tab:formalbase} shows the same lemmas on the baselines.
An unbound approval passes single use but fails both binding lemmas. A narrated card passes signature binding,
because the device does sign the digest, but fails effect binding and agent independence, because the card the
human reads is the agent's narration. A one-time code relayed through the agent's chat and an in-channel ``yes'' fail
effect binding and independence. Removing consumption from \sys{} breaks only single use, and a shared-key message authentication code (MAC) in
place of the passkey breaks only non-repudiation, which is the property aiAuthZ's symmetric receipts gave up
by design~\citep{kodathala2026aiauthz}.

\begin{table}[t]
\centering\small
\caption{The same lemmas on the baselines. \ok{} verified, \no{} falsified, n/a not stated for that model. V3 is a
mutant of \sys{} without consumption and V5 replaces the passkey with a shared-key MAC; neither is implemented as a
mode.}
\label{tab:formalbase}
\setlength{\tabcolsep}{5pt}
\begin{tabular}{@{}l c c c c c@{}}
\toprule
Model & L1 effect & L1 sig. & L2 single use & L3 indep. & L4 non-rep.\\
\midrule
\inputtab{tables/tab_formal_baselines.tex}
\bottomrule
\end{tabular}
\end{table}

\paragraph{What the model does not cover.} Free-text arguments that a grant leaves free are not modeled, and
\code{L5\_grant\_scope} makes no claim about them. Nonces and expiry are not modeled because the fresh approval
identifier already does their work in the model. The model also says nothing about whether a person reads the
card, which is \Hyp{1} and the subject of \cref{sec:eval}.
